\documentclass[dvipdfmx]{jsarticle}
\usepackage{tikz,ascmac,amsmath,amssymb,amsthm,bm,physics,arydshln,mathcomp,wrapfig,pgfplots,multicol}
\usetikzlibrary{calc}
\usetikzlibrary{intersections}
\DeclareMathOperator{\id}{id} %\idは定義されていない
\newtheorem*{answer*}{解答}

\begin{document}

・作りたい問題

・中間レポートで出た極値問題（井口微積p66参照）

・マクローリン展開したあとの不等式$\sin x > x-x^3/6$とか

・ラグランジュの未定乗数法と最大最小の関係について

・リーマン積分の定義計算について（ダルブー使わずに）

・合成関数の積分可能性

・コーシーアダマールってlimsupだけど、lim$a_n$が存在しないときのようなコーシーアダマールを使わざるを得ない状況

・一様収束性と$limsup$の関係、具体的な求め方

・ヘッセ行列の$(1,1)$成分が$0$の場合の極値

・マクローリン展開できない（解析的でない）関数

・不定積分できるが原始関数が存在しない関数

\end{document}
